% Paper 4, §14 — Contradictions live and open frontiers
% Thesis: method-pluralism is principled because the contradictions live at
% the boundaries between methods, and resolving them requires *cross-method*
% evidence — not better single-method tools. This section concentrates the
% [CONTRADICTION] markers raised across §3–§13 and, for each, names the
% strongest form of each side, the evidence on the table, and the experiment
% that would resolve it.

\section{Contradictions live and open frontiers}
\label{sec:contradictions}

The five methods of this paper agree often enough to make the field
move (the GPT-2-small IOI \glsterm{circuit}{circuit} is the worked convergence example;
the GPT-2-XL factual-recall locus is the second), but the cases where
they \emph{disagree} are not embarrassments to be papered over: they
are where the field's next experiments live. We collect five anchor
contradictions raised across \S\S\ref{sec:logit-lens-family}--\ref{sec:scaling-snapshot}
and discuss each in the same three-part form: the strongest version of
each side's claim, the evidence on the table as of 2026, and the
experiment that would resolve it. The closing argument recasts
method-pluralism as a principled stance: contradictions live at the
\emph{boundaries between} methods, and cross-method evidence is the
only currency that can spend them.

\subsection{The SAE canonical-feature claim}
\label{sec:contradictions-sae}

The strongest form of the \citet{bricken2023-monosemanticity} and
\citet{templeton2024-scaling-monosemanticity} program is that there
exists a privileged decomposition of activations into ``true features'' ---
a \glsterm{feature}{feature} dictionary that any sufficiently-careful \glsterm{sparse-autoencoder}{SAE} training
recovers, modulo a known equivalence. The strongest opposing form,
from \citet[\S 1]{leask2025-sae-not-canonical}, is that no such
canonical decomposition exists: meta-SAEs trained on the decoder
$\mathbf{W}_{\text{dec}}$ of a target \glsterm{sparse-autoencoder}{SAE} recover a sparse decomposition
of the target's atoms (the ``Einstein'' latent splits into
``scientist'', ``Germany'', ``famous person'') --- non-atomicity --- and
\glsterm{sae-stitching}{SAE stitching} reveals novel latents present in larger SAEs that smaller
SAEs cannot recover --- incompleteness.

The evidence on the table, summarized in \S\ref{sec:sae-canonicality}:
two ICLR-2025 lines of experiment supporting the negative; a corpus of
qualitative interpretability successes from
\citet{templeton2024-scaling-monosemanticity} consistent with the
positive but not entailing it; and the field's pragmatic resting
position --- ``apply SAEs of several widths''
\citep[\S 1]{leask2025-sae-not-canonical} --- which is methodological
agnosticism, not resolution.

\textbf{Resolving experiment.} A basis-finding criterion under which two
independent \glsterm{sparse-autoencoder}{SAE} training runs --- at different widths $M_1, M_2$,
different sparsity targets, different random seeds, on the same
activations --- converge on the same \glsterm{feature}{feature} inventory modulo a known
equivalence (rotation, sign, permutation, scale). Until that criterion
is exhibited, the canonical-\glsterm{feature}{feature} claim is unfalsifiable in the
direction that matters and the field is right to treat SAEs as useful
basis-dependent tools.

\subsection{Lens cross-model unreliability}
\label{sec:contradictions-lens}

The strongest form of the lens program
\citep{belrose2023-tuned-lens}: an affine translator
$(A_\ell, \mathbf{b}_\ell)$ trained by KL distillation against the
final-layer output is sufficient to absorb the basis drift that breaks
the bare \glsterm{logit-lens}{logit lens}, and the resulting tuned-lens trajectory is a
useful diagnostic on any pre-norm \glsterm{transformer}{transformer}. The strongest opposing
form, raised in \S\ref{sec:lens-frontier} and \S\ref{sec:scaling-no}:
the lens benchmarks predate the architectural substrate of the deployed
2026 lineup --- \glsterm{multi-head-latent-attention}{MLA}-compressed-KV (\glsterm{precision-pinning}{DeepSeek-V3}), fine-grained MoE,
\glsterm{nsa}{NSA}-sparse attention, sliding-window layers, hybrid attention/\glsterm{state-space-model}{SSM}
stacks --- and whether the affine-translator ansatz suffices on these
substrates is empirical, not deduced.

The evidence on the table is asymmetric. The \emph{positive} case for
extension is that the \glsterm{residual-stream}{residual stream} remains additive across these
variants; the \emph{negative} case, raised in
\S\ref{sec:scaling-no}, is the
\citet{halawi2023-overthinking} layer-of-prediction-vs-decision result
\textemdash\ early-layer predictions can be \emph{more} robust than the
final layer on incorrect-demonstration prompts, and choosing the right
layer has no a-priori recipe. \citet{chuang2023-dola}'s \glsterm{dola}{DoLa} contrast
turns the inconsistency into a decoding strategy on
LLaMA-7B/13B/33B, but whether \glsterm{dola}{DoLa}'s gains survive at frontier-2026
scale is not systematically reproduced.

\textbf{Resolving experiment.} Train tuned lenses across the
frontier-2026 lineup (Llama 3 / 4, Qwen 2.5 / 3, \glsterm{precision-pinning}{DeepSeek-V3}, Gemma 3),
report the per-architecture distillation-loss curves, and test whether
the trajectory diagnostics (smoothness, layer-of-acquisition,
contrastive \glsterm{dola}{DoLa} decoding) carry over with the architecture-specific
modifications they require. The empirical answer settles whether the
lens framework is a property of the additive \glsterm{residual-stream}{residual stream} (it
extends) or of a specific GPT-2-era attention/\glsterm{ffn}{FFN} profile (it does
not).

\subsection{Probe correlation versus causation}
\label{sec:contradictions-probe}

The strongest pre-2023 form of the probing program: high-accuracy
\glsterm{linear-probe}{linear probe} accuracy on a held-out test split implies the model
linearly represents the property at that layer. The strongest opposing
form, consolidated in
\citet[\S abstract]{belinkov2022-probing-survey} and the
\S\ref{sec:correlation-causation} contradiction block: \emph{decodable
does not entail used}. The model may compute a property at some
intermediate layer and discard it; later layers may use a different
representation; the \glsterm{probe-2}{probe} direction may correlate near-perfectly with
the used representation while not being it.

The evidence on the table is now substantial on both sides. The 2023+
direction-ablation literature
\citep[\S 3]{marks-tegmark-2023-truth} closes the gap on canonical
worked examples: ablating $\hat{\mathbf{d}}_{\text{MM}}$ at LLaMA-2-13B
mid-layers flips TRUE/FALSE behavior, validating the \glsterm{truth-direction-2}{truth direction}'s
\emph{use}. \deepencite{vennemeyer2025-sycophancy-not-one-thing §3,
multi-component-ablation} extends the pattern to a behavioral target
(\glsterm{sycophancy}{sycophancy}) with a multi-component decomposition. What remains open
is the alignment-relevant question: probes validated on
in-distribution evaluation prompts may not flag the same direction the
model uses on adversarial out-of-distribution inputs --- and that
generalization gap is precisely the load-bearing question for
\textbf{[Paper~5]} alignment-detection probing.

\textbf{Resolving experiment.} Pair every alignment-claim \glsterm{probe-2}{probe} with
direction-ablation on adversarial OOD evaluation suites; report
whether the ablation-induced behavior change persists under
distribution shift. A \glsterm{probe-2}{probe} whose ablation effect halves on a
held-out OOD shift has not validated the use claim it asserted.

\subsection{The truth direction across datasets}
\label{sec:contradictions-truth}

The strongest version of the linear-truth-representation program:
LLaMA-family models on factual statements admit a single linear axis
$\mathbf{d}_{\text{truth}} \in \R^D$, recoverable late-layer by
mass-mean and validated by patching
\citep[\S 3]{marks-tegmark-2023-truth}. The strongest opposing form
\citep[\S 4]{marks-tegmark-2023-truth} and the
\S\ref{sec:truth-frontier} contradiction block: per-dataset truth
directions are \emph{antipodal} early, rotate to orthogonal mid, and
align only late --- the ``stark misalignment'' of the original paper's
own §4. The 2025
\citet{vennemeyer2025-sycophancy-not-one-thing} \glsterm{sycophancy}{sycophancy}
decomposition generalizes the multi-component pattern: an
alignment-relevant phenomenon thought representationally singular
admits a three-direction (agreement / praise / genuine-agreement)
decomposition, each independently steerable. The evidence on the table
is therefore that ``behaviorally-singular properties may be
representationally-plural'' on at least two cases. CCS as the
unsupervised analog is documented to fail under negation
\citep[\S 1.1]{marks-tegmark-2023-truth} --- a near-truth confound, not
truth.

\textbf{Resolving experiment.} A multi-axis basis-recovery procedure
that, applied independently to \texttt{cities}, \texttt{neg\_cities},
\texttt{larger\_than}, and \texttt{cities\_conj} at matched layers,
returns a low-dimensional subspace whose per-dataset projections
recover the single-dataset directions modulo a known equivalence
(rotation, sign, scale). Until the subspace structure is exhibited,
``the \glsterm{truth-direction-2}{truth direction}'' is a useful one-dimensional projection of a
richer subspace, in a fixed model and dataset family --- not a
basis-free property of the model.

\subsection{Circuit-tracing scalability past Anthropic-2025}
\label{sec:contradictions-circuit-tracing}

The strongest form of the attribution-graph program
\citep{lindsey2025-circuit-tracing,lindsey2025-biology}: replacing
MLPs with cross-layer transcoders enables end-to-end Jacobian
attribution graphs on \glsterm{claude}{Claude}-class models, with reported findings
including forward planning in poetry, multi-step factual chains
(``Texas'' $\to$ ``Austin'' for ``the capital of the state Dallas is
in''), refusal-\glsterm{circuit}{circuit} features, and cross-lingual \glsterm{feature}{feature} overlap.
The strongest opposing form, raised in \S\ref{sec:circuit-forward}:
the Haiku findings have not been independently reproduced as of 2026,
the $\sim\!20\%$-of-pretraining \glsterm{transcoder}{transcoder}-training cost is a real
adoption barrier outside frontier labs, and \glsterm{transcoder}{transcoder} canonicality
inherits the \citet{leask2025-sae-not-canonical,dunefsky2024-transcoders}
critique --- a different \glsterm{transcoder}{transcoder} training run may factor the same
MLP into a different sparse basis, and the stability of attribution
graphs across \glsterm{transcoder}{transcoder} seeds is open empirical.

The evidence on the table is single-source on the substantive claims;
the 2025-06 Anthropic open-release of the \glsterm{circuit}{circuit}-tracer tooling
(\S\ref{sec:scaling-pipelines}) is a corrective force, but a corrective
force is not a replication.

\textbf{Resolving experiment.} Two transcoders trained on the same
\glsterm{claude}{Claude}-class model with different random seeds, and an attribution
graph on the same prompt computed against each. The conclusions ---
which features mediate the behavior, which features feed which, the
shape of the graph past the input embedding --- should agree modulo a
known \glsterm{transcoder}{transcoder} equivalence. If they don't, attribution graphs are
seed-dependent in the sense that matters for \glsterm{circuit}{circuit} claims, and the
canonicality concern transfers from SAEs to circuits in a load-bearing
way.

\begin{contradiction}
\textbf{Method-pluralism is principled, not provisional.} The five
contradictions above share a structural \glsterm{feature}{feature}: each one is the
\emph{boundary between} two of the paper's methods. \glsterm{sparse-autoencoder}{SAE} canonicality
is a question that probing alone cannot answer (probes don't decompose
activations into a basis) and that \glsterm{activation-patching}{activation patching} alone cannot
answer (patching tests components, not \glsterm{feature}{feature} inventories). Lens
extension to frontier-2026 architectures is a question that lens-only
benchmarks can only partially settle --- the ground truth requires a
patching-validated model of what the lens \emph{should} read out at
each layer. \glsterm{probe-2}{Probe} correlation-vs-causation is, by 2026 consensus, a
question whose resolution \emph{requires} pairing probes with
intervention machinery from \S\ref{sec:patching}. The truth-direction
plurality is visible only when probes are run across datasets
\emph{and} subspace methods (\citet{vennemeyer2025-sycophancy-not-one-thing})
are layered on top. \glsterm{circuit}{Circuit}-tracing scalability is, finally, a question
whose closure requires both \glsterm{transcoder}{transcoder}-replication evidence (within
the \glsterm{sparse-autoencoder}{SAE}/\glsterm{transcoder}{transcoder} method-family) and downstream behavior-edit
validation (with patching-style interventions). At every contradiction
listed above, the resolving experiment is \emph{cross-method}: it
reads off two methods on the same artifact and tests whether they
agree where they should. \emph{Better single-method tools cannot
close these gaps}, because what the gaps measure is the disagreement
between the questions different methods answer.
\end{contradiction}

\subsection{What ``partially read'' means, restated}
\label{sec:contradictions-partially-read}

Paper~4's title --- \emph{the internal computation can be partially
read} --- is now precise. We can extract human-legible structure for
\emph{some} behaviors (IOI on GPT-2 small, factual recall on GPT-2 XL,
\glsterm{truth-direction-2}{truth direction} in LLaMA-2-13B late layers, abstract features on
\glsterm{claude}{Claude} 3 Sonnet), on \emph{some} models (open-source GPT-family +
Gemma 2 + LLaMA 2 dominate the literature; the frontier-2026 lineup is
sparser), with \emph{cross-method agreement} on \emph{some} claims (the
two convergence cases of \S\ref{sec:cross-method}). The five anchor
contradictions catalog the open frontiers in the same \glsterm{vocabulary}{vocabulary}:
where the methods disagree, what each side asserts, what evidence is
on the table, and what experiment would resolve it. The closing
discipline of the paper is that progress is measured in
\emph{contradictions retired by cross-method evidence}, not in
single-method headline results --- and that the field, in 2026, has
the methodological \glsterm{vocabulary}{vocabulary} to do that work.

\subsection{Forward link}
\label{sec:contradictions-forward}

The five contradictions above are the explicit interpretability-side
input to the alignment-evaluation discussion of \textbf{[Paper~5]}.
Three load-bearing claims of the alignment-claim literature ---
``the model has a deception \glsterm{feature}{feature},'' ``the model uses its \glsterm{chain-of-thought}{CoT}
faithfully,'' ``\glsterm{activation-steering}{activation steering} is a safe inference-time
intervention'' --- depend on a \glsterm{probe-2}{probe}-and-\glsterm{feature}{feature} substrate whose
canonicality and use validity is exactly what \S\S\ref{sec:contradictions-probe},
\ref{sec:contradictions-truth}, and \ref{sec:contradictions-sae} have
just declared open. Paper~5 picks up the question of how much weight
those claims can carry given the evidence on the table here.
